\providecommand{\Lat}{\Lambda}
\providecommand{\gBKM}{\mathfrak{g}}
\providecommand{\phioo}{\varphi_{0,1}}
\providecommand{\phimt}{\varphi_{-2,1}}

\section*{Wave 14 Q5/20: Fresh cross-read of Lorgat 2020 paper and slides}

\paragraph{Scope.} Fresh holistic read of
\texttt{raeez.lorgat.automorphic-corrections.pdf} (10pp) and
the slide transcript \texttt{/tmp/automorphic-slides.txt} (953 lines).
Target: structural claims not yet adversarially audited by Waves
11--13. Five unsurfaced items emerge.

\subsection*{Claim A: the elliptic-genus factor of $2$}

Paper p.~2, after Theorem~2: \emph{``the elliptic genus of a K3
surface is given by the weak Jacobi form $\phioo$.''} Slide 919
(Arithmetic Geometry of some CY3's, Theorem panel): \emph{``the
elliptic genus of a K3 surface is $2\phioo$.''} The factor-of-two
discrepancy is material: the weak Jacobi form $\phioo$ of weight
$0$ index $1$ has $q^0 r^0$ constant term $10$, while the K3
elliptic genus $\mathrm{Ell}(K3; \tau, z) = 2\phioo$ has
$q^0 r^0$ constant term $20 = 2\chi_\mathrm{top}(K3) - 24 = 24-4 = 20$
(Eichler--Zagier convention, normalised so $\mathrm{Ell}(K3; \tau, 0)
= \chi_\mathrm{top}(K3) = 24$ at $z = 0$). The correct statement
is the slide version: $\mathrm{Ell}(K3) = 2\phioo$. The paper
statement is off by $\chi_\mathrm{top}/12$.

\subsection*{Claim B: sign of the Oberdieck--Pixton constant $C$}

Paper p.~2, Theorem~2: $Z^X(q,t,p) = C/(\Delta_5)^2$. Slide 917:
$Z^{S\times E} = -C/(\Delta_5)^2$. The Oberdieck--Pixton
\emph{Igusa cusp form conjecture} (arXiv 2018, now theorem) states
the generating function equals $-1/\chi_{10}(-q,-p,t) = -1/(-\Delta_5)^2
\cdot (\ldots)$; the precise sign depends on the refined versus
unrefined convention. The $(-p)^n$ factor in the summand absorbs
one sign; the slide sign $-C$ is the Oberdieck--Pixton
conventional form, while the paper's $+C$ is a transcription
error at p.~2 line (Theorem~2 RHS).

\subsection*{Claim C: lattice-polarisation-dependent GBKM family
$\gBKM_L$}

Paper p.~2 Conjecture~1 is terse: ``all eight diagonal divisor
modular forms of Gritsenko--Cl\'ery arise as reciprocal square
roots of $Z^X_{L, h_M}$, and these are denominator functions of
GBKMs with root multiplicities specified by $g_N - h_M$-twisted
twined elliptic genera of K3.'' Slides 934--936 give a
\emph{strictly stronger} version: \emph{``there is a GBKM algebra
with correction $\gBKM \subset \gBKM_L$, depending on $L$, given
by a $(g_N, h_M)$-twisted elliptic genus of a K3-surface
$F_L^{(g_N, h_M)}$, a Jacobi form of weight 0, index 1.''} The
slide version introduces a \emph{family parametrised by
lattice-polarisations $L$} of K3, with an explicit pair of
twisted elliptic genera $F_L^{(g_N, h_M)}$ as Jacobi-form inputs.
This family structure is absent from the paper. Wave 14 should
treat $L \mapsto \gBKM_L$ as a \emph{functor from the
Nikulin-admissible lattice-polarisation groupoid to GBKM
superalgebras}, with $\gBKM_{\emptyset} = \gBKM_{\Delta_5}$ the
unpolarised case.

\subsection*{Claim D: the second Jacobi form $\phimt = \nu_{11}^2/\eta^6$}

Slide 940 introduces $B = \phimt = \nu_{11}^2/\eta^6$, the weak
Jacobi form of weight $-2$ index $1$, and writes the case
analysis $N \in \{1,2,3,5,7\}$, $N=4$, $N=6$, $N=8$ using the
pair $(A, B) = (\tfrac14 \phioo^2, \phimt)$. The paper mentions
only $\phioo$. The two-form basis $(\phioo, \phimt)$ is the
standard Eichler--Zagier basis of weak Jacobi forms of index $1$
(their Theorem~9.3): every weak Jacobi form of index $1$ and
even weight $k$ is a combination $\alpha(\tau) \phioo + \beta(\tau)
\phimt$ with $\alpha \in M_k(\mathrm{SL}_2(\mathbb{Z}))$,
$\beta \in M_{k+2}(\mathrm{SL}_2(\mathbb{Z}))$. The twisted
elliptic genera $F_L^{(g_N, h_M)}$ therefore lie in a
two-dimensional space over the ring of
$\mathrm{SL}_2(\mathbb{Z})$-level-$N$ modular forms, with basis
$(\phioo, \phimt)$; the case split by $N$ in slides 943--953 is
explicit CHL-orbifold decomposition into this basis.

\subsection*{Claim E: pfaffian signature $(3,3) \twoheadrightarrow
(3,2)$ restriction}

Paper p.~4, paragraph below $\wedge^2$ definition: the pfaffian
scalar product on $\Lat^4 \wedge \Lat^4$ is \emph{signature
$(3,3)$ even unimodular}; the orthogonal complement of
$q = e_1\wedge e_3 + e_2 \wedge e_4$ is $\Lat^{3,2}$ of signature
$(3,2)$. The drop in rank is by a single timelike direction:
$q$ itself is a $(1,1)$-hyperbolic vector. Wave 11--13 audits
treated $\Lat^{3,2}$ as primitive; the \emph{embedding}
$\Lat^{3,2} \hookrightarrow \Lat^4 \wedge \Lat^4 \simeq
\Lat^{3,3}$ as $q^\perp$ is the arithmetic-geometric fact that
lifts $\mathrm{Sp}_4(\mathbb{Z}) \twoheadrightarrow \mathrm{O}(\Lat^{3,2})_+$
to a subquotient of $\mathrm{O}(\Lat^{3,3})$. Since $\Lat^{3,3}
\simeq \Lat^{2,2} \oplus \Lat^{1,1} \simeq U \oplus U \oplus U$
is the genus-$2$ Igusa lattice, this situates $\gBKM_{\Delta_5}$
inside the $\mathrm{O}(3,3)$-Borcherds programme. Relevance:
$\Lat^{3,3}$ carries $\mathrm{O}_3(3,3) = \mathrm{SL}_4 \times
\mathrm{SL}_4 / \mathbb{Z}_2$ exceptional isogenies, suggesting
a $\mathrm{GL}_4$-trace presentation of $\kappa_{\mathrm{BKM}}(\Delta_5) = 5$
via the $e_1 \wedge e_3 + e_2 \wedge e_4$ canonical form.

\subsection*{Synthesis for Wave 14}

The paper is a local construction of $\gBKM_{\Delta_5}$; the
slides generalise to the lattice-polarised family
$\{\gBKM_L\}_{L \in \mathcal{N}}$ indexed by Nikulin-admissible
polarisations $L$ with twisted elliptic genus $F_L^{(g_N, h_M)}$
as input. The pair $(A, B) = (\tfrac14 \phioo^2, \phimt)$ is the
Eichler--Zagier-basis in which $F_L^{(g_N, h_M)}$ expands.
Waves 11--13 audited $\kappa_{\mathrm{BKM}}(\Phi_N) = c_N(0)/2$ for
$N \in \{1,2,3,4,6\}$ using $\phioo$ alone; incorporating
$\phimt$ doubles the test surface, gives a second numerical
witness per case, and isolates the $\eta^{-6}$ factor as the
$\mathrm{SU}(2)_1$-WZW correction (level-$1$ character
denominator).

\paragraph{Four statements for Wave 14 synthesis.}
(i) $\mathrm{Ell}(K3) = 2\phioo$ (slide 919; paper p.~2 drops
the factor).
(ii) $Z^{S\times E} = -C/\Delta_5^2$ with explicit sign
(slide 917; paper p.~2 drops the sign).
(iii) Lattice-polarised GBKM family $L \mapsto \gBKM_L$ with
$\gBKM \subset \gBKM_L$ (slides 934--936; paper's
Conjecture~1 is the $L = \emptyset$ specialisation).
(iv) Eichler--Zagier basis $(\phioo, \phimt)$ with
$\phimt = \nu_{11}^2/\eta^6$ carries the $N=4, 6, 8$ CHL cases
(slides 940--953; paper uses only $\phioo$).

All four sharpen claims already on the table; none overturn
$\kappa_{\mathrm{BKM}}(\Phi_N) = c_N(0)/2$.
